% ══════════════════════════════════════════════════════════════════════════════
%  RL Primer: от нуля до PER — теоретическое введение
%  Для читателя, знакомого с ML, но не знакомого с RL.
%  Build: pdflatex rl_primer.tex  (дважды)
% ══════════════════════════════════════════════════════════════════════════════
\documentclass[a4paper, 11pt]{article}

\usepackage[utf8]{inputenc}
\usepackage[T2A]{fontenc}
\usepackage[english, russian]{babel}
\usepackage{amsmath, amssymb, amsthm}
\usepackage{graphicx}
\usepackage{geometry}
\geometry{left=2.5cm, right=2.5cm, top=2.5cm, bottom=2.8cm}
\usepackage[hidelinks,unicode]{hyperref}
\usepackage{xcolor}
\usepackage{tikz}
\usetikzlibrary{arrows.meta, positioning, backgrounds, calc}
\usepackage{booktabs}
\usepackage{enumitem}
\usepackage[framemethod=tikz]{mdframed}
\usepackage{microtype}
\usepackage{parskip}
\setlength{\parskip}{6pt}

% ── Цвета ────────────────────────────────────────────────────────────────────
\definecolor{defblue}{HTML}{1A3A6E}
\definecolor{boxbg}{HTML}{EDF2FB}
\definecolor{notebg}{HTML}{FEF9EC}
\definecolor{noteborder}{HTML}{E8A020}
\definecolor{codegray}{HTML}{F4F4F4}
\definecolor{mainblue}{HTML}{2255AA}

% ── Окружения ─────────────────────────────────────────────────────────────────
\mdfdefinestyle{defstyle}{%
  linecolor=defblue, linewidth=1.5pt,
  backgroundcolor=boxbg,
  innertopmargin=8pt, innerbottommargin=8pt,
  innerleftmargin=10pt, innerrightmargin=10pt,
  skipabove=8pt, skipbelow=8pt
}
\mdfdefinestyle{notestyle}{%
  linecolor=noteborder, linewidth=1pt,
  backgroundcolor=notebg,
  innertopmargin=7pt, innerbottommargin=7pt,
  innerleftmargin=10pt, innerrightmargin=10pt,
  skipabove=6pt, skipbelow=6pt
}

\newenvironment{defbox}[1][]{%
  \begin{mdframed}[style=defstyle]
  \ifx&#1&\else{\bfseries\color{defblue}#1.}\quad\fi\ignorespaces
}{%
  \end{mdframed}
}

\newenvironment{notebox}{%
  \begin{mdframed}[style=notestyle]
  {\bfseries\color{noteborder}Интуиция.}\quad\ignorespaces
}{%
  \end{mdframed}
}

% ── Команды ───────────────────────────────────────────────────────────────────
\newcommand{\E}{\mathbb{E}}
\newcommand{\R}{\mathbb{R}}
\DeclareMathOperator*{\argmax}{arg\,max}

% ── Заголовок ─────────────────────────────────────────────────────────────────
\title{%
  \textbf{Обучение с подкреплением: от основ до PER}\\[6pt]
  \large Теоретическое введение для понимания работы\\
  \normalsize\color{gray}«Зависимость качества DQN-агента self-play
  от параметра приоритизации $\alpha$ в игре «Уголки»»%
}
\author{}
\date{}

% ══════════════════════════════════════════════════════════════════════════════
\begin{document}
\maketitle
\vspace{-1em}
{\color{gray}\hrule}
\vspace{0.5em}
\tableofcontents
\newpage

% ══════════════════════════════════════════════════════════════════════════════
\section{Зачем RL — и чем он отличается от обычного ML}
% ══════════════════════════════════════════════════════════════════════════════

В supervised learning у нас есть датасет $\{(x_i, y_i)\}$: вход и правильный
ответ. Задача — научить модель отображению $x \mapsto y$.

В обучении с подкреплением \textbf{правильного ответа нет}. Есть \textbf{агент},
который совершает действия в \textbf{среде}, и среда в ответ даёт
\textbf{скалярную награду} (reward). Цель агента — выработать стратегию
поведения, которая максимизирует суммарную награду.

\begin{notebox}
Аналогия из жизни: обучение вождению. Инструктор не говорит «сейчас поверни
руль на 23 градуса». Он говорит «хорошо» или «плохо», и вы постепенно сами
нащупываете правильное поведение. RL — это примерно то же самое.
\end{notebox}

Ключевые отличия от supervised learning:

\begin{enumerate}[itemsep=2pt]
  \item \textbf{Нет разметки.} Агент не получает «правильное действие» —
        он получает только награду, которая может прийти много ходов спустя.
  \item \textbf{Отложенное вознаграждение.} Действие сейчас влияет на
        ситуацию через много шагов. Это называется \emph{temporal credit
        assignment problem}.
  \item \textbf{Распределение данных зависит от агента.} В отличие от
        фиксированного датасета, какие ситуации мы будем видеть —
        зависит от того, как агент себя ведёт.
\end{enumerate}

% ══════════════════════════════════════════════════════════════════════════════
\section{Марковский процесс принятия решений (MDP)}
% ══════════════════════════════════════════════════════════════════════════════

Формальная математическая основа RL — \textbf{Марковский процесс принятия
решений} (Markov Decision Process, MDP).

\begin{defbox}[MDP]
MDP задаётся кортежем $(\mathcal{S}, \mathcal{A}, P, R, \gamma)$, где:
\begin{itemize}[itemsep=2pt]
  \item $\mathcal{S}$ — пространство состояний,
  \item $\mathcal{A}$ — пространство действий,
  \item $P(s' \mid s, a)$ — функция переходов: вероятность оказаться в
        состоянии $s'$, если в состоянии $s$ выбрать действие $a$,
  \item $R(s, a, s')$ — немедленная награда,
  \item $\gamma \in [0, 1)$ — коэффициент дисконтирования.
\end{itemize}
\end{defbox}

\textbf{Марковское свойство} означает, что следующее состояние $s'$ зависит
только от текущего состояния $s$ и действия $a$, но не от всей предыстории.
Это принципиально: мы можем строить оптимальные стратегии, смотря только на
текущее состояние.

\subsection{Как это выглядит на практике}

На каждом шаге $t$:
\begin{enumerate}[itemsep=2pt]
  \item Агент наблюдает состояние $s_t$.
  \item Агент выбирает действие $a_t$.
  \item Среда переходит в новое состояние $s_{t+1} \sim P(\cdot \mid s_t, a_t)$.
  \item Агент получает награду $r_t = R(s_t, a_t, s_{t+1})$.
\end{enumerate}

\begin{center}
\begin{tikzpicture}[
  font=\small,
  box/.style={draw=defblue!60, fill=boxbg, rounded corners=6pt,
              minimum width=2.4cm, minimum height=1.0cm, align=center, thick},
  arr/.style={-{Stealth[length=6pt]}, thick, defblue!70}
]
  \node[box] (ag) {Агент};
  \node[box, right=3.5cm of ag] (env) {Среда};

  \draw[arr] (ag.north) -- ++(0, 0.5) -| node[above, pos=0.25] {действие $a_t$} (env.north);
  \draw[arr] (env.south) -- ++(0,-0.5) -| node[below, pos=0.25] {состояние $s_{t+1}$, \; награда $r_t$} (ag.south);
\end{tikzpicture}
\end{center}

\subsection{Пример из нашей работы}

В игре «Уголки»:
\begin{itemize}[itemsep=2pt]
  \item $\mathcal{S}$ — все возможные позиции фишек на доске $8 \times 8$;
        закодировано как тензор $3 \times 8 \times 8$.
  \item $\mathcal{A}$ — все пары клеток (откуда, куда): $64 \times 64 = 4096$
        действий. Большинство нелегальны.
  \item $P$ — детерминированный (игра без случайности).
  \item $R$ — $+100$ за победу, $-100$ за поражение, $+0{,}2 \times \Delta d$
        за приближение к цели, $-0{,}01$ за каждый ход.
  \item $\gamma = 0{,}99$.
\end{itemize}

% ══════════════════════════════════════════════════════════════════════════════
\section{Политика и функции ценности}
% ══════════════════════════════════════════════════════════════════════════════

\subsection{Политика}

\begin{defbox}[Политика]
Политика $\pi$ — это стратегия агента. Стохастическая политика задаётся
распределением $\pi(a \mid s)$ — вероятностью выбрать действие $a$ в
состоянии $s$. Детерминированная политика задаётся функцией $a = \pi(s)$.
\end{defbox}

\subsection{Суммарная дисконтированная отдача}

После момента времени $t$ агент получит следующие награды: $r_t, r_{t+1},
r_{t+2}, \ldots$ Мы хотим максимизировать не одну награду, а всю
\textbf{будущую отдачу}:

\begin{equation}
  G_t = r_t + \gamma\, r_{t+1} + \gamma^2 r_{t+2} + \cdots
       = \sum_{k=0}^{\infty} \gamma^k r_{t+k}
  \label{eq:return}
\end{equation}

Зачем нужен $\gamma < 1$? Во-первых, математически: сумма сходится.
Во-вторых, содержательно: награда сейчас ценнее, чем та же награда потом
(как в экономике). Типичное значение $\gamma = 0{,}99$ означает, что
награда через 100 шагов стоит $0{,}99^{100} \approx 0{,}37$ от её
номинала — это «горизонт» примерно 100 шагов.

\subsection{Функция ценности состояния}

\begin{defbox}[Функция $V^\pi$]
Ценность состояния $s$ при политике $\pi$ — ожидаемая будущая отдача,
если начать из $s$ и следовать $\pi$:
\[
  V^\pi(s) = \E_\pi\bigl[G_t \;\big|\; s_t = s\bigr]
\]
\end{defbox}

$V^\pi$ отвечает на вопрос: «Насколько хорошо оказаться в состоянии $s$,
если дальше следовать политике $\pi$?»

\subsection{Функция ценности действия (Q-функция)}

\begin{defbox}[Q-функция]
Ценность действия $a$ в состоянии $s$ при политике $\pi$:
\[
  Q^\pi(s, a) = \E_\pi\bigl[G_t \;\big|\; s_t = s,\, a_t = a\bigr]
\]
\end{defbox}

$Q^\pi(s,a)$ отвечает на вопрос: «Насколько хорошо выбрать именно $a$ в $s$,
а потом следовать $\pi$?»

Связь между ними:
\[
  V^\pi(s) = \sum_a \pi(a \mid s)\, Q^\pi(s, a)
\]

\subsection{Оптимальные функции ценности}

Нас интересует \textbf{оптимальная} политика $\pi^*$ — та, которая
максимизирует $V^\pi(s)$ для всех состояний одновременно. Оптимальные
функции ценности:
\[
  V^*(s) = \max_\pi V^\pi(s), \qquad
  Q^*(s, a) = \max_\pi Q^\pi(s, a)
\]

Связь между ними фундаментальна: если мы знаем $Q^*$, оптимальная
политика тривиальна:
\[
  \pi^*(s) = \argmax_{a} Q^*(s, a)
\]

Именно поэтому Q-обучение так популярно: достаточно научить сеть
приближать $Q^*$.

% ══════════════════════════════════════════════════════════════════════════════
\section{Уравнения Беллмана}
% ══════════════════════════════════════════════════════════════════════════════

Уравнения Беллмана — это рекуррентные соотношения, которые связывают
ценность текущего состояния с ценностями следующих. Они лежат в основе
всех алгоритмов RL.

\subsection{Уравнение Беллмана для $Q^\pi$}

Раскроем определение $Q^\pi$ через определение $G_t$:

\begin{equation}
  Q^\pi(s, a) = \E\bigl[r + \gamma\, Q^\pi(s', \pi(s')) \;\big|\; s, a\bigr]
  \label{eq:bellman_policy}
\end{equation}

Иными словами, ценность $Q^\pi(s,a)$ равна немедленной награде плюс
дисконтированная ценность в следующем состоянии при продолжении
следования $\pi$.

\subsection{Оптимальное уравнение Беллмана}

Для оптимальной $Q^*$:

\begin{equation}
  \boxed{Q^*(s, a) = \E\bigl[r + \gamma \max_{a'} Q^*(s', a')
                    \;\big|\; s, a\bigr]}
  \label{eq:bellman_opt}
\end{equation}

Это уравнение Беллмана оптимальности. Оно самосогласованное: $Q^*$ и
есть решение этого уравнения. Именно к нему будет сходиться Q-обучение.

\begin{notebox}
Аналогия с ML. В supervised learning у нас есть функция потерь
$\mathcal{L}(\theta) = \E[(y - f_\theta(x))^2]$ и мы ищем $\theta^*$
минимизирующий её. В RL задача хитрее: «правильный ответ» $r + \gamma
\max_{a'} Q^*(s', a')$ сам зависит от $Q^*$, которую мы ищем.
Поэтому нужны специальные методы.
\end{notebox}

% ══════════════════════════════════════════════════════════════════════════════
\section{Табличное Q-обучение}
% ══════════════════════════════════════════════════════════════════════════════

Самый простой способ найти $Q^*$ — хранить все значения $Q(s, a)$ в
большой таблице и постепенно обновлять их.

\subsection{TD-ошибка}

На каждом шаге мы наблюдаем переход $(s, a, r, s')$ и вычисляем
\textbf{ошибку временной разности} (temporal difference error, TD-ошибка):

\begin{equation}
  \delta = \underbrace{r + \gamma \max_{a'} Q(s', a')}_{\text{TD-цель}}
           - \underbrace{Q(s, a)}_{\text{текущая оценка}}
  \label{eq:td_error}
\end{equation}

Затем обновляем таблицу:
\[
  Q(s, a) \leftarrow Q(s, a) + \alpha_{\rm lr}\, \delta
\]

где $\alpha_{\rm lr}$ — шаг обучения (не путать с параметром PER $\alpha$!).

\begin{notebox}
Аналогия с SGD. TD-обновление — это как один шаг градиентного спуска по
функции потерь $\mathcal{L} = \delta^2 / 2$, только «правильный ответ»
(TD-цель) временно фиксируется. Это и есть суть bootstrapping в RL.
\end{notebox}

\subsection{Сходимость и ограничения}

Алгоритм Q-обучения был предложен Уоткинсом в его докторской диссертации
(1989) и формально проанализирован совместно с Дэяном в работе
\cite{Watkins92}. Авторы доказали, что при двух условиях — каждая пара
$(s, a)$ посещается бесконечно часто и шаг обучения $\alpha_{\rm lr}$
убывает по правилу Робинса--Монро — Q-обучение сходится к $Q^*$ с
вероятностью 1. Это фундаментальный теоретический результат, на котором
стоит весь Q-learning и его потомки.

Проблема: если $|\mathcal{S}| \times |\mathcal{A}|$ велико
(в нашем случае: $\sim 10^{30} \times 4096$), таблица не помещается
в память и условие «посетить каждую пару бесконечно часто» невыполнимо
на практике. Нужна аппроксимация функции $Q$.

% ══════════════════════════════════════════════════════════════════════════════
\section{Deep Q-Network (DQN)}
% ══════════════════════════════════════════════════════════════════════════════

\subsection{Нейросеть вместо таблицы}

Прорыв совершила работа Мниха и соавторов из DeepMind \cite{Mnih15},
опубликованная в \textit{Nature} в 2015~году. Авторы показали, что
\textbf{одна нейронная сеть} способна играть в 49 различных видеоигр
Atari на уровне человека или выше, получая на вход только пиксели экрана и
счёт — без какого-либо специфичного для игры знания. До этого считалось,
что каждая задача требует своей ручной настройки. Ключевых нововведений
было два: Experience Replay и Target Network — оба описаны ниже.

Идея DQN \cite{Mnih15}: заменить таблицу нейронной сетью
$Q_\theta(s, a)$ с параметрами $\theta$. Функция потерь:

\begin{equation}
  \mathcal{L}(\theta) = \E_{(s,a,r,s') \sim \mathcal{D}}
    \Bigl[\bigl(\underbrace{r + \gamma \max_{a'} Q_{\theta^-}(s', a')}_{y
    \text{ — TD-цель}} - Q_\theta(s, a)\bigr)^2\Bigr]
  \label{eq:dqn_loss}
\end{equation}

Обратите внимание на $\theta^-$ в TD-цели — об этом ниже.

В нашей работе:
\[
  Q_\theta: \mathbb{R}^{3 \times 8 \times 8} \to \mathbb{R}^{4096}
\]
Сеть принимает на вход тензор состояния и выдаёт \textbf{сразу все}
$Q$-значения для всех $4096$ действий за один прямой проход.

\subsection{Проблема 1: корреляция в онлайн-обучении}
\label{sec:replay}

Если обучать сеть сразу на каждом переходе $(s_t, a_t, r_t, s_{t+1})$ в
порядке их получения, данные будут сильно скоррелированы: соседние
состояния похожи, градиенты «тянут» сеть в одном направлении, и обучение
нестабильно.

\textbf{Решение: Experience Replay} \cite{Lin92}.

Идея буфера опыта была предложена Линем ещё в 1992~году в контексте
обучения роботов: вместо того чтобы выбрасывать каждый переход после
одного использования, его сохраняют и многократно «переигрывают».
Лин показал, что это существенно повышает эффективность использования
данных в нескольких задачах управления. Мних и соавторы \cite{Mnih15}
заново открыли и популяризировали эту идею применительно к глубоким
нейросетям, превратив её в стандартный компонент DQN-алгоритмов.

\begin{defbox}[Буфер опыта]
Все переходы $(s, a, r, s', \text{done})$ записываются в кольцевой буфер
$\mathcal{D}$ размером $C$ (у нас $C = 10^5$). На каждом шаге из буфера
случайно сэмплируется мини-батч размером $B$ и делается один шаг
градиентного спуска.
\end{defbox}

Это аналог перемешивания датасета в supervised learning: разрываем
корреляции, получаем i.i.d.-подобные батчи.

\begin{notebox}
Буфер работает как «память» агента. Старые переходы постепенно
вытесняются новыми по мере заполнения. Размер буфера — компромисс:
маленький буфер означает мало разнообразия, слишком большой — много
устаревших данных от старой политики.
\end{notebox}

\subsection{Проблема 2: нестабильность TD-цели}

Если TD-цель $y = r + \gamma \max_{a'} Q_\theta(s', a')$ вычисляется той
же сетью $\theta$, которую мы обновляем, то цель «убегает» во время
обучения — мы гонимся за постоянно движущейся мишенью. Это приводит к
расходимости.

\textbf{Решение: Target Network} (целевая сеть) \cite{Mnih15}. Это второе
ключевое нововведение оригинальной статьи DeepMind.

\begin{defbox}[Target Network]
Ведём две копии сети: \textbf{online}-сеть $Q_\theta$ (обучается каждый шаг)
и \textbf{target}-сеть $Q_{\theta^-}$ (заморожена). TD-цель считается по
$Q_{\theta^-}$. Каждые $T_{\rm target}$ шагов делаем жёсткое копирование:
$\theta^- \leftarrow \theta$.

У нас $T_{\rm target} = 1000$ шагов.
\end{defbox}

\subsection{Проблема 3: исследование среды}

Если агент всегда выбирает действие с максимальным $Q$, он может
застрять в локальном оптимуме. Нужен баланс между \emph{exploitation}
(использовать знания) и \emph{exploration} (исследовать неизвестное).

\textbf{Решение: $\varepsilon$-жадная политика.}

\begin{defbox}[$\varepsilon$-жадная политика]
\[
  a_t = \begin{cases}
    \text{случайное легальное действие} & \text{с вероятностью } \varepsilon \\
    \argmax_{a} Q_\theta(s_t, a)       & \text{с вероятностью } 1 - \varepsilon
  \end{cases}
\]
$\varepsilon$ линейно убывает от $1{,}0$ до $0{,}05$ за 50\,000 шагов:
сначала почти чистое исследование, потом почти чистая эксплуатация.
\end{defbox}

\subsection{Маскировка нелегальных действий}

В нашей задаче многие из 4096 действий нелегальны (фишки нет на клетке,
нельзя двигаться так). Перед выбором действия нелегальные маскируются:
\[
  Q_\theta(s, a) \leftarrow -\infty \quad \text{для всех нелегальных } a
\]
Тогда $\argmax$ гарантированно вернёт легальное действие.

\subsection{Итоговый алгоритм DQN}

\begin{enumerate}[itemsep=3pt]
  \item Инициализировать $\theta$, $\theta^- \leftarrow \theta$, буфер $\mathcal{D}$.
  \item Для каждого шага:
    \begin{enumerate}[itemsep=2pt]
      \item Наблюдать $s_t$, выбрать $a_t$ по $\varepsilon$-жадной политике.
      \item Сделать ход, получить $r_t, s_{t+1}$, записать в $\mathcal{D}$.
      \item Сэмплировать батч из $\mathcal{D}$.
      \item Вычислить TD-цели: $y_i = r_i + \gamma \max_{a'} Q_{\theta^-}(s'_i, a')$.
      \item Сделать шаг градиентного спуска: $\theta \leftarrow \theta -
            \nabla_\theta \frac{1}{B}\sum_i (y_i - Q_\theta(s_i, a_i))^2$.
      \item Обновить $\theta^-$ каждые $T_{\rm target}$ шагов.
    \end{enumerate}
\end{enumerate}

% ══════════════════════════════════════════════════════════════════════════════
\section{Self-Play}
% ══════════════════════════════════════════════════════════════════════════════

\subsection{Идея}

В однопользовательских задачах RL (например, Atari) есть агент и
фиксированная среда. В играх с нулевой суммой (как «Уголки») у нас
два игрока. Self-play — это режим, в котором \textbf{агент играет сам с
собой}: обе стороны используют одну и ту же сеть.

Это изящное решение: по мере того как сеть улучшается, соперник тоже
становится сильнее, создавая постоянное давление на дальнейший рост.

Исторически self-play хорошо зарекомендовал себя в нардах (TD-Gammon,
Tesauro, 1994) и достиг пика известности в серии работ AlphaGo/AlphaZero
компании DeepMind: в~2017~году AlphaZero обучилась игре в шахматы, сёги и
го \textbf{только через self-play}, не используя ни одной партии человека,
и превзошла все специализированные программы. Мы адаптируем ту же идею к
«Уголкам»: вместо огромной вычислительной инфраструктуры AlphaZero —
лёгкий DQN, который, тем не менее, успешно осваивает игровую стратегию.

\subsection{Одна сеть для обоих игроков}

Ключевая техника — \textbf{канонический фрейм}. Доска всегда
кодируется с точки зрения текущего игрока: его фишки — на своей позиции,
чужие — на чужой. Для Игрока $-1$ доска поворачивается на $180°$.

Таким образом, сеть всегда видит одно и то же: «свои фишки сверху-слева,
надо перенести их вниз-вправо» — независимо от того, кто ходит.

\subsection{Negamax bootstrapping}

В обычном DQN TD-цель:
\[
  y = r + \gamma \max_{a'} Q_{\theta^-}(s', a')
\]
Но в self-play после хода Игрока 1 ход переходит к Игроку $-1$. Состояние
$s'$ кодируется уже \emph{с точки зрения Игрока $-1$}. А игра нулевая
сумма, значит:
\[
  V(s \text{ для Игрока 1}) = -V(s \text{ для Игрока } {-1})
\]

Поэтому TD-цель в self-play корректируется знаком:

\begin{equation}
  \boxed{y = r - \gamma \max_{a'} Q_{\theta^-}(s', a')}
  \label{eq:negamax}
\end{equation}

Это называется \textbf{negamax bootstrapping}. Без этой поправки цели
были бы неверными и сеть обучалась бы «помогать» противнику.

\subsection{Нестационарность}

В обычном RL среда фиксирована — $P(s'|s,a)$ не меняется. В self-play
«соперник» — это та же сеть, которая обновляется. С точки зрения каждого
из игроков среда \textbf{нестационарна}: распределение $P$ меняется
вместе с параметрами $\theta$.

Это важно для понимания того, почему гиперпараметры из статических сред
(например, Atari) могут не работать в self-play.

% ══════════════════════════════════════════════════════════════════════════════
\section{Prioritized Experience Replay (PER)}
% ══════════════════════════════════════════════════════════════════════════════

\subsection{Мотивация}

Метод Prioritized Experience Replay был предложен Шаулем и соавторами из
DeepMind \cite{Schaul16} и представлен на ICLR~2016. Авторы показали, что
на~\textbf{41 из 49 игр Atari} агент с PER опережает стандартный DQN с
равномерным буфером при одинаковом числе шагов обучения. В оригинальной
работе в качестве рабочего значения параметра приоритизации выбрано
$\alpha = 0{,}6$ — оно оказалось лучшим компромиссом между
полной приоритизацией и равномерностью именно для статичных сред Atari.
Вопрос о том, сохраняется ли этот результат в нестационарных средах,
таких как self-play, и является предметом нашей работы.

В стандартном Experience Replay все переходы сэмплируются равномерно.
Но не все переходы одинаково полезны:

\begin{itemize}[itemsep=3pt]
  \item Переход с большой TD-ошибкой $|\delta|$ означает, что сеть
        сильно ошиблась в своей оценке — это «неожиданный» и потенциально
        информативный переход.
  \item Переход с малой $|\delta|$ означает, что сеть уже хорошо знает
        это место — повторное обучение на нём принесёт мало пользы.
\end{itemize}

\begin{notebox}
Аналогия из обучения: если вы хорошо решаете задачи типа A и плохо типа B,
разумно чаще практиковать B. PER делает то же самое — «сложные» переходы
сэмплируются чаще.
\end{notebox}

\subsection{Формула приоритетов}

\begin{defbox}[Приоритеты в PER \cite{Schaul16}]
Каждому переходу $i$ в буфере назначается приоритет:
\[
  p_i = |\delta_i| + \varepsilon
\]
где $\varepsilon = 10^{-6}$ — маленькая константа, чтобы никакой переход
не получил нулевой приоритет. Вероятность сэмплирования:
\[
  P(i) = \frac{p_i^{\,\alpha}}{\displaystyle\sum_{j} p_j^{\,\alpha}}
\]
\end{defbox}

\subsection{Параметр $\alpha$ — главный предмет исследования}

Параметр $\alpha \in [0, 1]$ контролирует \textbf{степень приоритизации}:

\begin{center}
\begin{tabular}{cl}
  \toprule
  Значение $\alpha$ & Поведение \\
  \midrule
  $\alpha = 0$ & $P(i) = \text{const}$ — равномерная выборка (стандартный DQN) \\
  $\alpha = 0{,}6$ & умеренная приоритизация (стандарт для Atari) \\
  $\alpha = 1$ & $P(i) \propto |\delta_i|$ — полная приоритизация по TD-ошибке \\
  \bottomrule
\end{tabular}
\end{center}

\begin{figure}[h]
\centering
\begin{tikzpicture}[font=\small, scale=1.0]
  \draw[-{Stealth}, thick, gray] (-0.2,0) -- (4.5,0)
    node[right] {$|\delta_i|$};
  \draw[-{Stealth}, thick, gray] (0,-0.2) -- (0,3.2)
    node[above] {$P(i)$};

  % alpha=0: flat
  \draw[gray, thick] (0.2, 0.85) -- (4.1, 0.85);
  \node[right, gray] at (4.1, 0.85) {\small $\alpha=0$};

  % alpha=0.6
  \draw[mainblue, thick]
    (0.2, 0.45) .. controls (1.5, 0.9) and (2.8, 1.8) .. (4.1, 2.5);
  \node[right, mainblue] at (4.1, 2.5) {\small $\alpha=0{,}6$};

  % alpha=1
  \draw[defblue!80, thick, dashed]
    (0.2, 0.15) .. controls (0.9, 0.4) and (2.3, 1.7) .. (4.1, 3.05);
  \node[right, defblue!80] at (4.1, 3.05) {\small $\alpha=1$};

  % tick labels
  \foreach \x/\l in {1/малая, 2/средняя, 3/большая}
    \node[below, font=\scriptsize, gray] at (\x, 0) {\l};
  \node[below, font=\scriptsize, gray] at (2, -0.45) {TD-ошибка $|\delta_i|$};
\end{tikzpicture}
\caption{Зависимость вероятности выборки от TD-ошибки при разных $\alpha$.
При $\alpha=0$ все переходы равновероятны; при $\alpha=1$ переходы с
большой ошибкой доминируют.}
\end{figure}

\subsection{Проблема смещения (bias) и IS-коррекция}

Когда мы сэмплируем не равномерно, оценка градиента становится
смещённой: переходы с высоким приоритетом «перегружают» функцию потерь.
Чтобы это скомпенсировать, каждому переходу назначается
\textbf{importance sampling (IS) вес}:

\begin{equation}
  w_i = \frac{1}{N \cdot P(i)}\,^{-\beta}
      = \bigl(N \cdot P(i)\bigr)^{-\beta}
  \label{eq:is_weight_raw}
\end{equation}

Для численной стабильности нормируем на $\max_j w_j$:
\[
  \tilde{w}_i = \frac{w_i}{\max_j w_j}
\]

Скорректированная функция потерь:
\begin{equation}
  \mathcal{L}_{\rm PER}(\theta)
  = \frac{1}{B} \sum_{i=1}^{B}
    \tilde{w}_i \cdot \bigl(y_i - Q_\theta(s_i, a_i)\bigr)^2
  \label{eq:per_loss}
\end{equation}

\subsection{Параметр $\beta$ — степень коррекции}

$\beta \in [0, 1]$ контролирует степень IS-коррекции:
\begin{itemize}[itemsep=2pt]
  \item $\beta = 0$: IS-веса все равны 1 — коррекции нет.
  \item $\beta = 1$: полная коррекция смещения.
\end{itemize}

На практике $\beta$ \textbf{анилируется} от $0{,}4$ до $1{,}0$ по мере
обучения. В начале обучения агрессивная коррекция мешает — сеть ещё
нестабильна. К концу, когда $Q$-значения уже более точные, нужна полная
коррекция для гарантий сходимости.

В нашей работе: $\beta: 0{,}4 \to 1{,}0$ за 500\,000 шагов.

\subsection{Полный алгоритм PER}

По сравнению с обычным DQN добавляются три операции:

\begin{enumerate}[itemsep=3pt]
  \item \textbf{При записи} нового перехода в буфер: назначить ему
        максимальный текущий приоритет (гарантирует, что новый переход
        будет отобран хотя бы раз).
  \item \textbf{При сэмплировании}: сэмплировать пропорционально
        $P(i)$, вычислить IS-веса $\tilde{w}_i$.
  \item \textbf{После gradient step}: пересчитать $\delta_i$ для
        отобранных переходов и \textbf{обновить} их приоритеты в буфере.
\end{enumerate}

% ══════════════════════════════════════════════════════════════════════════════
\section{Всё вместе: что и зачем в нашей работе}
% ══════════════════════════════════════════════════════════════════════════════

Теперь можно понять каждый компонент.

\subsection{Полный pipeline}

\begin{enumerate}[itemsep=5pt]
  \item \textbf{Среда} — ортогональные «Уголки» $8 \times 8$. MDP
        с детерминированными переходами и reward shaping.

  \item \textbf{Кодирование} — тензор $3 \times 8 \times 8$ в
        каноническом фрейме. Один тензор описывает состояние с точки
        зрения любого из двух игроков.

  \item \textbf{DQN-сеть} — 2 свёрточных слоя + 2 линейных,
        выход — 4096 Q-значений. Нелегальные действия маскируются
        $-\infty$ перед $\argmax$.

  \item \textbf{Self-play} — оба игрока используют одну сеть.
        Negamax bootstrapping \eqref{eq:negamax} обеспечивает
        правильный знак TD-цели в игре с нулевой суммой.

  \item \textbf{Imitation warm-start} — перед self-play агент
        несколько итераций обучается поведенческому клонированию
        от эвристического агента. Это даёт разумную стартовую
        политику и ускоряет дальнейшее обучение.

  \item \textbf{PER (или Uniform)} — буфер опыта $10^5$ переходов.
        В варианте PER сэмплирование пропорционально $(|\delta_i| +
        \varepsilon)^\alpha$; приоритеты обновляются после каждого
        gradient step.
\end{enumerate}

\subsection{Что именно исследуется}

Работа ставит один вопрос: \textbf{как параметр $\alpha$ влияет на
конечное качество политики} при фиксированных всех остальных
гиперпараметрах?

Проверяются четыре значения: $\alpha \in \{0,\; 0{,}3,\; 0{,}6,\; 0{,}9\}$.
$\alpha = 0$ соответствует стандартному DQN без приоритизации.

\subsection{Гипотеза о нестационарности}

Стандартное значение $\alpha = 0{,}6$ было подобрано для задач Atari, где
среда фиксирована. В self-play «соперник» — это та же обновляемая сеть.
Это означает:

\begin{itemize}[itemsep=3pt]
  \item TD-ошибки прошлых переходов быстро устаревают: то, что было
        «трудным» для старой политики, может стать лёгким для новой.
  \item При высоком $\alpha$ такие устаревшие «трудные» переходы
        продолжают доминировать в выборке → перекошенный градиент.
  \item IS-коррекция ($\beta$) компенсирует смещение, но в начале
        обучения $\beta = 0{,}4$ и коррекция неполная.
\end{itemize}

Гипотеза: \textbf{агрессивная приоритизация дестабилизирует self-play},
приводя к пиковым TD-ошибкам и снижению качества политики.

Дополнительный аргумент в пользу осторожности при выборе $\alpha$ даёт
работа Хесселя и соавторов \cite{Hessel18}, создавших \textbf{Rainbow} —
комбинацию шести улучшений DQN одновременно (включая PER, Double DQN,
Dueling и другие). Авторы обнаружили, что гиперпараметры компонентов
\emph{взаимодействуют}: оптимальное $\alpha$ в изоляции может оказаться
неоптимальным в сочетании с другими методами. В нашем случае
«дополнительным компонентом», меняющим характер обучения, выступает
именно self-play.

Таким образом, литература даёт основания полагать, что $\alpha = 0{,}6$
не является универсально лучшим выбором и может требовать адаптации под
конкретную среду.

\subsection{Метрика и оценка}

Агент оценивается в 1000 партиях против Greedy и Heuristic агентов.
Метрика — \textbf{доля побед среди результативных партий}:
\[
  \text{win rate} = \frac{\text{победы}}{\text{победы} + \text{поражения}}
\]
Ничьи исключены, потому что при лимите 300 ходов их очень много
(57–74\% против Random) — они заглушают сигнал.

Для статистики используется \textbf{bootstrap confidence interval}:
из результатов $n$ партий делается 10\,000 ресэмплов с возвращением,
по каждому считается win rate, берётся 2.5-й и 97.5-й перцентили.
Это даёт 95\% ДИ, не требуя нормальности распределения.

% ══════════════════════════════════════════════════════════════════════════════
\section*{Краткий словарь}
% ══════════════════════════════════════════════════════════════════════════════

\begin{tabular}{@{}p{3.8cm}p{10cm}@{}}
  \toprule
  Термин & Пояснение \\
  \midrule
  Agent / Агент & Обучаемая система, выбирающая действия. \\
  Environment / Среда & Всё, с чем взаимодействует агент. \\
  State $s$ & Полное описание текущей ситуации. \\
  Action $a$ & Ход, который делает агент. \\
  Reward $r$ & Скалярный сигнал обратной связи от среды. \\
  Return $G_t$ & Дисконтированная сумма будущих наград \eqref{eq:return}. \\
  $\gamma$ & Коэффициент дисконтирования (у нас 0.99). \\
  Policy $\pi$ & Стратегия: как выбирать действие в состоянии. \\
  $Q(s,a)$ & Ожидаемая отдача при действии $a$ в $s$. \\
  TD-ошибка $\delta$ & Разность TD-цели и текущей оценки $Q$ \eqref{eq:td_error}. \\
  Experience Replay & Буфер переходов для разрыва корреляций. \\
  Target Network & Замороженная копия сети для стабильных TD-целей. \\
  $\varepsilon$-жадная & Политика: случайное действие с вероятностью $\varepsilon$. \\
  Self-play & Агент играет против собственной копии. \\
  Negamax & Знаковая поправка TD-цели для zero-sum игр \eqref{eq:negamax}. \\
  PER & Prioritized Experience Replay — выборка по $|\delta|$ \cite{Schaul16}. \\
  $\alpha$ (PER) & Степень приоритизации. $\alpha=0$ — равномерно. \\
  $\beta$ (PER) & Степень IS-коррекции. $\beta=1$ — полная коррекция. \\
  \bottomrule
\end{tabular}

% ══════════════════════════════════════════════════════════════════════════════
\begin{thebibliography}{9}
\bibitem{Watkins92}
  \textit{Watkins C.J.C.H., Dayan P.}
  Q-learning.
  \textit{Machine Learning}, 8:279--292, 1992.

\bibitem{Mnih15}
  \textit{Mnih V. et~al.}
  Human-level control through deep reinforcement learning.
  \textit{Nature}, 518:529--533, 2015.

\bibitem{Schaul16}
  \textit{Schaul T. et~al.}
  Prioritized Experience Replay.
  \textit{ICLR}, 2016.

\bibitem{Lin92}
  \textit{Lin L.-J.}
  Self-improving reactive agents based on reinforcement learning,
  planning and teaching.
  \textit{Machine Learning}, 8:293--321, 1992.

\bibitem{Hessel18}
  \textit{Hessel M. et~al.}
  Rainbow: Combining improvements in deep reinforcement learning.
  \textit{AAAI}, 2018.
\end{thebibliography}

\end{document}
